\begin{thebibliography}{99}
\small
\raggedright
\setlength{\itemsep}{0pt}
\bibitem{erdos1948} P.~Erd\H{o}s,
  \href{https://users.renyi.hu/~p_erdos/1948-04.pdf}{\emph{On arithmetical
  properties of Lambert series}},
  J. Indian Math. Soc. (N.S.) \textbf{12} (1948), 63--66.
\bibitem{erdos1988} P.~Erd\H{o}s, \emph{On the irrationality of certain series:
  problems and results}, in A.~Baker (ed.), \emph{New Advances in Transcendence
  Theory}, Cambridge UP, 1988, pp.~102--109,
  doi:\href{https://doi.org/10.1017/CBO9780511897184.009}{10.1017/CBO9780511897184.009}.
\bibitem{borwein1991} P.~B. Borwein, \emph{On the irrationality of
  $\sum1/(q^{n}+r)$},
  J. Number Theory \textbf{37} (1991), no.~3, 253--259,
  doi:\href{https://doi.org/10.1016/S0022-314X(05)80041-1}{10.1016/S0022-314X(05)80041-1}.
\bibitem{borwein1992} P.~B. Borwein, \emph{On the irrationality of certain
  series}, Math. Proc. Cambridge Philos. Soc. \textbf{112} (1992), no.~1,
  141--146,
  doi:\href{https://doi.org/10.1017/S030500410007081X}{10.1017/S030500410007081X}.
  Van Assche cites Lemma~2 for the neighbouring little-$q$-Legendre
  evaluation used above.
\bibitem{az1998} T.~Amdeberhan and D.~Zeilberger, \emph{$q$-Ap\'ery
  irrationality proofs by $q$-WZ pairs}, Adv. Appl. Math. \textbf{20} (1998),
  no.~2, 275--283,
  \href{https://arxiv.org/abs/math/9804122}{arXiv:math/9804122},
  doi:\href{https://doi.org/10.1006/aama.1997.0565}{10.1006/aama.1997.0565}.
\bibitem{bv1994} P.~Bundschuh and K.~V\"a\"an\"anen,
  \href{https://numdam.org/item/CM_1994__91_2_175_0.pdf}{\emph{Arithmetical
  investigations of a certain infinite product}},
  Compositio Math. \textbf{91} (1994), no.~2, 175--199.
\bibitem{zudilin2004} W.~Zudilin,
  \href{https://geodesic.mathdoc.fr/articles/10.4064/aa111-2-4/}{\emph{Heine's
  basic transform and a permutation group for \(q\)-harmonic series}},
  Acta Arith. \textbf{111} (2004), no.~2, 153--164,
  doi:10.4064/aa111-2-4.
\bibitem{zudilin2016} W.~Zudilin,
  \href{https://arxiv.org/abs/1601.02688}{\emph{On the irrationality of
  generalized \(q\)-logarithm}}, arXiv:1601.02688; Res. Number Theory
  \textbf{2} (2016),
  doi:\href{https://doi.org/10.1007/s40993-016-0042-x}{10.1007/s40993-016-0042-x}.
  The remark that the results extend to non-integer \(p=r/s\), \(|p|>1\), under
  an assumption \(\log|r|>c\log|s|\) for a computable \(c>0\), is at the end of
  Section~2; no value of \(c\) is computed there, and the remark is made for the
  generalized \(q\)-logarithm of that paper.
\bibitem{duverney1996} D.~Duverney,
  \href{https://numdam.org/item/JTNB_1996__8_1_173_0.pdf}{\emph{\`A propos de la
  s\'erie $\sum_{n\ge1}x^{n}/(q^{n}-1)$}},
  J. Th\'eor. Nombres Bordeaux \textbf{8} (1996), no.~1, 173--181.
  Th\'eor\`eme~2 on p.~174 gives the rational-base region
  $\log|s|/\log|r|<\frac13(1-3/\pi^{2})=0.2320\ldots$ for this series;
  Th\'eor\`eme~1, for a general numerator, is weaker.
\bibitem{matalaaho2006} T.~Matala-aho, K.~V\"a\"an\"anen and W.~Zudilin,
  \href{https://doi.org/10.1090/S0025-5718-05-01812-0}{\emph{New irrationality
  measures for \(q\)-logarithms}}, Math. Comp. \textbf{75} (2006), no.~254,
  879--889, doi:10.1090/S0025-5718-05-01812-0.  The hypothesis
  \(p=1/q\in\Z\mathbin{\backslash}\{0,\pm1\}\) is carried in the abstract on p.~879 and in
  both theorem statements on p.~880, where the authors also record that their
  methods do not sharpen the \(q\)-harmonic case of~\cite{zudilin2004}.
\bibitem{adamczewskibell2013} B.~Adamczewski and J.~P. Bell,
  \href{https://arxiv.org/abs/1303.2019}{\emph{A problem about Mahler
  functions}}, Ann. Sc. Norm. Super. Pisa Cl. Sci. \textbf{17} (2017), no.~4,
  1301--1355; arXiv:1303.2019, 2013.  Theorem~1: over a field of characteristic
  zero, a power series is both \(k\)- and \(\ell\)-Mahler for multiplicatively
  independent \(k,\ell\) if and only if it is a rational function.
\bibitem{bellsmertnig2026} J.~Bell and D.~Smertnig,
  \href{https://arxiv.org/abs/2603.23456}{\emph{Mahler series with
  multiplicative coefficient sequences}}, arXiv:2603.23456v1,
  24 March 2026.  The introduction explicitly includes the divisor function
  among the examples which are not \(k\)-Mahler for any \(k\ge2\).
\bibitem{rhinviola1996} G.~Rhin and C.~Viola,
  \href{https://geodesic.mathdoc.fr/articles/10.4064/aa-77-1-23-56/}
  {\emph{On a permutation group related to $\zeta(2)$}},
  Acta Arith. \textbf{77} (1996), no.~1, 23--56,
  doi:10.4064/aa-77-1-23-56.
\bibitem{vanassche2001} W.~Van Assche,
  \href{https://arxiv.org/abs/math/0101187}{\emph{Little \(q\)-Legendre
  polynomials and irrationality of certain Lambert series}},
  Ramanujan J. \textbf{5} (2001), no.~3, 295--310,
  doi:\href{https://doi.org/10.1023/A:1012930828917}{10.1023/A:1012930828917}.
\bibitem{vandehey2013} J.~Vandehey,
  \href{https://arxiv.org/abs/1206.0340}{\emph{On an incomplete argument of
  Erd\H{o}s on the irrationality of Lambert series}},
  Integers \textbf{13} (2013), Paper A58.
\bibitem{lucatachiya2017} F.~Luca and Y.~Tachiya,
  \href{https://www.kurims.kyoto-u.ac.jp/~kyodo/kokyuroku/contents/pdf/2014-14.pdf}
  {\emph{Linear independence results for the values of divisor functions
  series}}, RIMS K\^oky\^uroku No.~2014 (2017), 138--150.
  Theorem~A on p.~139 restates the periodic-coefficient irrationality theorem;
  Example~1 on p.~140 gives the divisor-function specialization.
\bibitem{rivin2026} I.~Rivin,
\href{https://arxiv.org/abs/2604.25151}{\emph{Zero Coefficients of Rational
  Power Series and Rational Lambert Series}}, arXiv:2604.25151v1, 2026.
  Theorem~1.1 is on p.~2 and proved on pp.~6--7; the periodic-coefficient
  Corollary~6.4 is on p.~9.
\bibitem{kovactao2024} V.~Kova\v{c} and T.~Tao,
  \href{https://arxiv.org/abs/2406.17593}{\emph{On several irrationality
  problems for Ahmes series}}, Acta Math. Hungar. \textbf{175} (2025),
  572--608; arXiv:2406.17593, 2024.
\bibitem{erdosproblems} T.~F. Bloom,
  \href{https://www.erdosproblems.com/1049}{\emph{Erd\H{o}s Problem \#1049}},
  \texttt{erdosproblems.com/1049}, accessed 28 July 2026
  (page displays ``last edited 28 September 2025'').  The current record
  labels the problem open, cites \textup{[Er88c, p.~102]} and \textup{[Er48]},
  and explicitly describes its status as the website owner's present
  assessment, with no guarantee of literature completeness.
\bibitem{formalconjectures1049} The Formal Conjectures Authors,
  \href{https://github.com/google-deepmind/formal-conjectures/blob/f776d2f2039351b00737ffcafb9d7d7666e1d9af/FormalConjectures/ErdosProblems/1049.lean}
  {\emph{FormalConjectures.ErdosProblems.\texttt{1049}}},
  Lean source at commit \texttt{f776d2f}, 2026, accessed 28 July 2026.  The
  irrationality declarations are unproved; the Lambert-series identity is
  proved.
\end{thebibliography}

\end{document}
